\documentclass[11pt]{article}
\usepackage[a4paper,margin=1in]{geometry}
\usepackage{amsmath}
\usepackage{amssymb}
\usepackage{booktabs}
\usepackage{hyperref}
\hypersetup{colorlinks=true, linkcolor=blue, citecolor=blue, urlcolor=blue}

\title{A Posture-Adaptive Belt for Reducing Adverse Sacroiliac Joint
Loading During Sitting:\\ Design Rationale and Simulation Study}

\author{Hermes Agent\thanks{Corresponding build agent, Hermes Agent by Nous
Research.} \and Walker Kirkpatrick, ND}
\date{October 17, 2026}

\begin{document}
\maketitle

\begin{abstract}
Sacroiliac joint (SIJ) pain is a recognized source of chronic low back
pain, and pelvic belts are a standard conservative management tool. We
identify a mechanical failure mode of \emph{static} belts: a belt sized to
hold trousers while standing becomes several times tighter once its
wearer sits, because lower-body circumferences increase from standing to
sitting and the pelvis rotates backward, redirecting belt force across
lax or previously injured SIJs. We present a lightweight, posture-adaptive
belt design---a six-axis inertial sensor, a debounced posture classifier,
a per-activity tension policy, a fail-loose safety layer, and a
rate-limited circumference actuator---implemented as a dependency-free
Python control stack and evaluated in closed-loop simulation against a
frozen static-belt baseline on identical sensor streams. Over a simulated
full day (standing, driving, and desk work), the adaptive belt reduced
the seated adverse (misdirected) SIJ load integral by
89.3\% (from 80,579 to 8,642~N\,s), kept
100.0\% of seated time at or under the comfort floor, and completed
the day with zero safety-rule violations on roughly 1.9~mAh of
battery. Seated slack is biomechanically defensible: pelvic-belt tension
beyond roughly 50~N does not further reduce SIJ laxity, so relaxing when
seated should cost nothing therapeutically. The design is a simulation
and control reference implementation, not a medical device.
\end{abstract}

\section{Introduction}
Sacroiliac joint pain is one of the most common causes of low back pain,
accounting for an estimated 15--30\% of cases in reviews of the
diagnostic literature \cite{ref1}, with SIJ dysfunction estimated as the
primary pain source in 10--25\% of affected patients \cite{ref2}.
Recognized etiologies include pregnancy, repetitive stress, and trauma;
posttraumatic arthritis of the SI joints can follow injuries such as
motor vehicle collisions, and SIJ dysfunction is associated with
osteoarthritis and the spondyloarthropathies \cite{ref3}. For a subset
of these patients---notably those with ligamentous laxity or prior
pelvic trauma---ordinary clothing is itself a daily hazard. A trouser
belt sized to hold up trousers while standing has a smaller
circumference than the pelvis presents when its owner sits down:
lower-body circumferences measurably increase with the switch from
standing to sitting \cite{ref11}, and the pelvis tends to move backward
from standing to sitting \cite{ref24}. The same belt, unchanged, then
applies a sharper and differently directed force across the SIJs at the
seated hip angle---precisely the situation this work addresses.

Belts are nonetheless the correct conservative tool. A pelvic belt worn
with a small force can generate a self-bracing effect at the SIJ
\cite{ref19}; application of a belt significantly reduces SIJ mobility
\cite{ref19}; pelvic belts may be beneficial in affected postpartum
patients \cite{ref3}; and belting is a recognized nonoperative option
alongside exercise, manual therapy, and orthotics \cite{ref10}. In a
five-day trial of strenuous manual labor, spine pain increased by only
0.2 on a numeric rating scale during the week participants used a
Serola sacroiliac belt, versus 0.9 without it \cite{ref8}.

Crucially, the dose--response of belts saturates: in Damen's study of
pelvic-belt positioning, a belt tension of 100~N did not reduce SIJ
laxity more than 50~N \cite{ref4}. Extra tension buys no extra
stabilization. This is the biomechanical license for the central design
move of this paper: a belt that \emph{relaxes} when its wearer sits
should lose nothing therapeutically while eliminating the seated
over-tension entirely.

A motorized consumer ``smart belt'' that self-adjusts its fit has
existed since at least 2015 \cite{ref16}; what consumer prior art lacks
is a clinical rationale, a posture-aware control policy, per-wearer
calibration, and a fail-loose safety architecture.

\section{Background}
\subsection{SIJ belt mechanics}
The therapeutic model is force closure: the belt adds an external
compression that supplements the muscle-generated force closure of the
pelvic ring \cite{ref19}. Two findings shape our design. First, the
effect saturates with tension \cite{ref4}. Second, positioning matters:
in Damen's experiment the belt was most effective in a high position
\cite{ref4}. Both argue for \emph{modulated} tension rather than
maximal tension---which is what a posture-adaptive actuator provides.

\subsection{Posture sensing from the waist}
Wearable technology is a promising alternative to laboratory systems
for movement and posture analysis \cite{ref13}. Waist-worn
accelerometers are the standard for physical-activity assessment, but
the sedentary-behavior literature cautions that waist-worn devices
alone distinguish sitting from light activity less cleanly than
thigh-worn inclinometry \cite{ref23}. Our design answers this caveat
explicitly rather than ignoring it: the classifier debounces posture
changes over a persistence window, gates switches on confidence, and
degrades to a loose command when it cannot name the posture. Confusing
seated driving with seated desk sitting is deliberately tolerated---both
receive slack targets, so the cost of that confusion is near zero.

\section{Methods}
\subsection{Design overview}
The belt is a closed loop: a six-axis inertial measurement unit (IMU)
at the iliac crest feeds a posture classifier; the classifier feeds a
per-activity tension policy; the policy feeds a fail-loose safety
monitor; and the safety monitor commands a rate-limited actuator that
changes the belt's own circumference. Every wearer-dependent number
lives in a JSON profile with a calibration wizard, so all geometry,
stiffness, timing, and tension targets are per-wearer customizable.

\subsection{First-order biomechanical model}
A belt can pull but never push, so belt tension follows a spring law
\begin{equation}
T \;=\; k \,\max\!\left(0,\; C_{\text{wearer}}(a) - B\right),
\end{equation}
where $C_{\text{wearer}}(a)$ is the wearer's pelvic circumference at
activity $a$ (standing, or standing plus a seated expansion
$\Delta C_{\text{sit}}$ when seated), $B$ is the belt circumference, and
$k$ is the belt's elastic stiffness. Projecting the belt force onto SIJ
axes at a posterior pelvic rotation $\theta$ gives a beneficial
force-closure component and an adverse, misdirected component:
\begin{align}
F_{\text{closure}} &= T\cos\theta,\\
F_{\text{adverse}} &= T\sin\theta .
\end{align}
We emphasize that this projection is a first-order engineering model:
it is used \emph{only} for relative comparisons between belt strategies
on the same simulated wearer, never as an absolute estimate of joint
loading. Posterior pelvic rotation is near zero standing and larger
sitting and driving, consistent with the observed backward pelvic
motion from standing to sitting \cite{ref24}.

Table~\ref{tab:assumptions} lists the model parameters, their default
values, and the provenance of each: cited direction where the
literature supports it, explicit engineering assumption where it does
not.

\begin{table}[t]
\centering
\caption{Model parameters, defaults, and provenance.}
\label{tab:assumptions}
\begin{tabular}{lll}
\toprule
Parameter & Default & Provenance \\
\midrule
Standing circumference & 880 mm & wearer calibration \\
Seated expansion $\Delta C_{\text{sit}}$ & 25 mm & increases when seated \cite{ref11}; \\
 & & magnitude is wearer-specific \\
Belt stiffness $k$ & 0.8 N/mm & engineering assumption \\
Standing hold tension & 8 N & engineering assumption \\
Seated floor tension & 3 N & design choice (slack) \\
Safety ceiling & 30 N & design limit \\
Rotation priors & 2/15/25$^\circ$ & direction cited \cite{ref24}; \\
 & & magnitudes assumed \\
Debounce window & 1.0 s & design choice (cf.\ \cite{ref23}) \\
Re-tighten delay & 2.0 s & design choice \\
Actuator slew & 40 mm/s & engineering assumption \\
\bottomrule
\end{tabular}
\end{table}

\subsection{Control policy and safety architecture}
The tension policy maps posture to intent: standing keeps a modest
trouser-holding tension; sitting and driving relax to a floor near
slack; \emph{unknown} means loose. Re-tightening is delayed by a grace
window after standing so the belt never cinches mid-motion; loosening
is never delayed. The safety layer enforces, in precedence order:
manual loosen override, manual freeze override, a sensor watchdog
(stale estimate $\rightarrow$ loose), a tension ceiling, and only then
the controller's request. The fail-safe state of the entire system is
the loosest state: since additional tension does not add stabilization
\cite{ref4}, a belt that fails loose cannot over-compress an SIJ, and
sagging trousers are a cosmetic failure rather than a clinical one.

\subsection{Simulation protocol}
The stack is exercised by a deterministic simulated IMU that plays a
posture scenario through the full pipeline at 50~Hz. Each experiment
runs the \emph{same} sensor stream (identical seed) through two belts:

\begin{enumerate}
\item the \textbf{adaptive belt}, the full stack as designed; and
\item the \textbf{static baseline}, an actuator frozen at the
circumference that yields the standing hold tension---i.e., ordinary
clothing, sized standing and left alone.
\end{enumerate}

Table~\ref{tab:day} shows the full-day scenario. The headline metric is
the seated adverse-load integral $\int F_{\text{adverse}}\,dt$ over
seated ticks---exactly the quantity a lax or previously injured SIJ
should not receive while its owner sits at a desk or in a car.

\begin{table}[t]
\centering
\caption{Simulated day schedule.}
\label{tab:day}
\begin{tabular}{clrrr}
\toprule
\# & Posture & Start (s) & End (s) & Minutes \\
\midrule
1 & Standing & 0 & 900 & 15 \\
2 & Driving & 900 & 2100 & 20 \\
3 & Standing & 2100 & 2160 & 1 \\
4 & Sitting & 2160 & 5760 & 60 \\
5 & Standing & 5760 & 6060 & 5 \\
6 & Sitting & 6060 & 9660 & 60 \\
7 & Driving & 9660 & 10860 & 20 \\
8 & Standing & 10860 & 11760 & 15 \\
\bottomrule
\end{tabular}
\end{table}

\section{Results}
\subsection{Full simulated day}
Over the 196-minute day (588,000 control ticks), the adaptive
belt reduced the seated adverse-load integral from
80,579~N\,s to 8,642~N\,s---a 89.3\%
reduction---while the static belt's seated tension peaked at
28.0~N against 27.4~N for the adaptive belt
(Table~\ref{tab:results}). The adaptive belt kept
100.0\% of seated ticks at or under the comfort floor, performed
8 actuator adjustments, used approximately 1.9~mAh of
battery (roughly 154~days on a 300~mAh cell), and completed
the day with zero safety-rule violations.

\begin{table}[t]
\centering
\caption{Seated adverse SIJ load: adaptive vs.\ static belt (full day,
seed 7).}
\label{tab:results}
\begin{tabular}{lrr}
\toprule
Metric & Adaptive & Static \\
\midrule
Adverse-load integral (N\,s) & 8,642 & 80,579 \\
Peak seated adverse force (N) & 27.4ADV & 28.0ADV \\
Peak seated tension (N) & 27.4 & 28.0 \\
Seated time at/under floor (\%) & 100.0 & --- \\
Actuator adjustments & 8 & 0 \\
Battery used (mAh) & 1.9 & --- \\
Safety events & 0 & 0 \\
\bottomrule
\end{tabular}
\end{table}

\subsection{Commute scenario}
The minimal version of the motivating problem---stand, drive for half
an hour, stand---reproduces the same pattern over a
50-minute commute: 89.3\% reduction in seated
adverse load (21,299 $\rightarrow$ 2,286~N\,s), with
zero safety events.

\section{Discussion}
Three design principles carry the result. First, \emph{seated slack is
free}: because belt tension beyond roughly 50~N does not further
reduce SIJ laxity \cite{ref4}, the therapeutic cost of relaxing when
seated is approximately zero, while the seated over-tension of a static
belt is the exact harm mode the design removes. Second,
\emph{fail-loose everywhere}: every fault path---dead sensor, stale
estimate, low confidence, ceiling breach---terminates at the loosest
belt position, and the boot state is loose. Third,
\emph{customize rather than assume}: every geometry, stiffness, timing,
and tension number is a profile field with a calibration wizard, so the
system adapts to the wearer rather than the reverse.

\section{Limitations}
The SIJ load model is a first-order projection used only for relative
comparisons; it is not validated against instrumented joint loading,
and absolute newton values should not be read clinically. The simulated
sensor is an idealized waist-worn IMU; the sedentary-behavior
literature cautions that waist-worn sensing alone distinguishes
postures less cleanly than thigh-worn sensing \cite{ref23}, and
although the classifier answers that caveat with debouncing, confidence
gates, and graceful degradation, hardware validation is future work.
Sitting-versus-driving confusion is tolerated by design and would need
revisiting if the two postures ever received different targets. The
evaluation is simulation only, with no human subjects, and no claim of
clinical efficacy is made; any physical implementation would require
appropriate review, and use for a diagnosed SIJ condition should be
supervised by a clinician.

\section{Conclusion}
A belt that senses its wearer's posture and relaxes when they sit
eliminates the seated over-tension that static belts impose on people
with sacroiliac joint pain, at approximately zero biomechanical cost.
The reference implementation---a dependency-free Python control stack
with a CLI, a reproducible simulation harness, and a 127-test suite---
reduces seated adverse SIJ loading by 89.3\% over a simulated
day without a single safety-rule violation. The design is released
open-source as a foundation for hardware prototyping and, eventually,
properly supervised clinical evaluation.

\section*{Reproducibility}
All results in this paper are produced by
\texttt{sibelt simulate --scenario day --json} and
\texttt{sibelt simulate --scenario commute --json} from the released
repository (seed 7); the raw reports are committed alongside this
paper in \texttt{docs/paper/results/}.

\bibliographystyle{plain}
\begin{thebibliography}{99}
\bibitem{ref1} Successful Diagnosis of Sacroiliac Joint Dysfunction (PMC). Accessed 2026-10-17. \url{https://pmc.ncbi.nlm.nih.gov/articles/PMC8517984}

\bibitem{ref2} Association Between SIJ Dysfunction and LBP (PMC). Accessed 2026-10-17. \url{https://pmc.ncbi.nlm.nih.gov/articles/PMC8043903}

\bibitem{ref3} SI Joint Dysfunction: Diagnosis and Treatment (AAFP). Accessed 2026-10-17. \url{https://www.aafp.org/afp/2022/0300/p239}

\bibitem{ref4} Damen: Does a pelvic belt influence sacroiliac joint laxity?. Accessed 2026-10-17. \url{https://repub.eur.nl/pub/72524}

\bibitem{ref8} Serola SI Belt: pain and functional impairment (PMC). Accessed 2026-10-17. \url{https://pmc.ncbi.nlm.nih.gov/articles/PMC9479175}

\bibitem{ref10} Nonoperative Treatment Options for SIJ Pain (PMC). Accessed 2026-10-17. \url{https://pmc.ncbi.nlm.nih.gov/articles/PMC7041669}

\bibitem{ref11} Kinematic 3D Human-Body Model: standing vs sitting garment fit (PMC). Accessed 2026-10-17. \url{https://www.ncbi.nlm.nih.gov/pmc/articles/PMC8464815}

\bibitem{ref13} Commercial Postural Devices: A Review (PMC). Accessed 2026-10-17. \url{https://pmc.ncbi.nlm.nih.gov/articles/PMC6929158}

\bibitem{ref16} Belty self-adjusting smart belt (Gizmodo). Accessed 2026-10-17. \url{https://gizmodo.com/the-only-gadget-the-world-really-needs-is-a-self-adjust-1677432880}

\bibitem{ref19} POGP: pelvic support belts for perinatal PGP. Accessed 2026-10-17. \url{https://thepogp.co.uk/_userfiles/pages/files/gps_msb_pgp_06.pdf}

\bibitem{ref23} Measuring sedentary behavior: waist- vs thigh-worn monitors. Accessed 2026-10-17. \url{https://d-nb.info/1254931295/34}

\bibitem{ref24} Influence of Pelvic Tilt on Femoral Anteversion (supine/standing/sitting). Accessed 2026-10-17. \url{https://onlinelibrary.wiley.com/doi/10.1111/os.12543}

\end{thebibliography}

\end{document}
